\section*{Hoofdstuk \ref{ch:b3:quantitative-geomorphology} --- Kwantitatieve geomorfologie}

\begin{solution}{exo:b3:quantitative-geomorphology:1}
$a/\sqrt{\pi\kappa t} = 0.30 - 0.04 = 0.26$, dus $\kappa t = (2/0.26)^2/\pi =
\qty{18.8}{m^2}$, en $t = 18.8/(5\times10^{-3}) = \qty{3.8}{kyr}$, binnen een
factor twee in beide richtingen, gezien de spreiding van $\kappa$.
\end{solution}

\begin{solution}{exo:b3:quantitative-geomorphology:2}
$E = \kappa\,|\nabla^2 z| = 4\times10^{-3}\times0.012 =
\qty{4.8e-5}{m/yr}$, \qty{0.048}{mm/yr}. Reliëf $E L^2/(2\kappa) =
0.012\times80^2/2 = \qty{38.4}{m}$.
\end{solution}

\begin{solution}{exo:b3:quantitative-geomorphology:3}
$E = 2\times10^{-6}\times(25\times10^6)^{0.5}\times0.01 = 2\times10^{-6}
\times5000\times0.01 = \qty{1.0e-4}{m/yr}$, \qty{0.1}{mm/yr}. Omdat $E$
in \unit{m/yr} staat, $A^m$ in $\unit{m}^{2m}$ en $S$ geen eenheid heeft, ligt $K$ in
$\unit{m}^{1-2m}\,\unit{yr^{-1}}$: \unit{yr^{-1}} voor $m = 0.5$.
\end{solution}

\begin{solution}{exo:b3:quantitative-geomorphology:4}
(a) Vast: een harde rotsbodem steilt de rivier op waar hij aan het oppervlak komt, en de trap
blijft bij de rotsbodem. (b) Verankerd aan de breuk: zolang de breuk verschuift,
de rivier wordt steiler waar hij die kruist. (c) Mobiel: de opheffing van de
kust verlaagde in één keer het \omterm{def:b3:quantitative-geomorphology:knickpoint}{basisniveau} van elke rivier; de \omterm{def:b3:quantitative-geomorphology:knickpoint}{knikpunten}
migreren stroomopwaarts en liggen, zoals de kinematische-golfstelling voorspelt voor $n
= 1$, op één hoogtelijn.
\end{solution}

\begin{solution}{exo:b3:quantitative-geomorphology:5}
$\theta = m/n = 0.5$. $S = (U/K)^{1/2}A^{-1/2} = (10^{-4}/10^{-9})^{1/2}
\times(10^7)^{-1/2} = 316\times3.16\times10^{-4} = 0.10$. Het verdubbelen van
$U$ vermenigvuldigt $k_s = (U/K)^{1/n}$ met $\sqrt2 = 1.41$; voor $n = 1$ met 2: hoe
groter $n$, hoe minder de steilheid op opheffing reageert.
\end{solution}

\begin{solution}{exo:b3:quantitative-geomorphology:6}
\sloppy Met $x$ stroomopwaarts vanaf de uitmonding, $\mathrm{d}x = -\mathrm{d}d$, dus $\chi(d)
= \int_d^L (A_0/c\,d'^{h})^{\theta}\,\mathrm{d}d' = A_0^{\theta} c^{-\theta}
\int_d^L d'^{-h\theta}\,\mathrm{d}d' = A_0^{\theta}c^{-\theta}(L^p -
d^p)/p$. Hier is $p = 0.19$, $A_0^{\theta}c^{-\theta} = 10^{2.7}/2^{0.45}$,
en $\chi = \qty{2935}{m}$.\par
\end{solution}

\begin{solution}{exo:b3:quantitative-geomorphology:7}
$c = 2\times10^{-6}\times(10^8)^{0.5} = \qty{0.02}{m/yr}$, \qty{20}{m/kyr};
\qty{5}{km} met die snelheid kost \qty{250}{kyr}. Het oppervlak, en daarmee de
celeriteit neemt stroomopwaarts af, dus het \omterm{def:b3:quantitative-geomorphology:knickpoint}{knikpunt} vertraagt terwijl het omhoog klimt.
\end{solution}

\begin{solution}{exo:b3:quantitative-geomorphology:8}
$\bar e = 0.40\times600 = \qty{240}{m}$; de toppen stijgen met $0.80\times240 =
\qty{192}{m}$; de gemiddelde hoogte daalt met $240 - 192 = \qty{48}{m}$.
\end{solution}

\begin{solution}{exo:b3:quantitative-geomorphology:9}
$h = h_0\ln(P_0/E)$: $0.5\ln 6 = \qty{0.90}{m}$ en $0.5\ln1.5 =
\qty{0.20}{m}$. Bij \qty{0.4}{mm/yr} overtreft de erosie de snelste productie,
$P_0$: de bodem wordt weggeschuurd en de helling wordt kaal gesteente, waarvan
de daling door verwering wordt beperkt.
\end{solution}

\begin{solution}{exo:b3:quantitative-geomorphology:10}
In de \omterm{def:b3:quantitative-geomorphology:steady-profile}{stationaire toestand} geldt $\mathrm{d}q/\mathrm{d}x = E$ en $q = 0$ bij de waterscheiding, dus
$q = E x$. Dan is $\kappa S = E x\,(1 - S^2/S_c^2)$, een kwadratische vergelijking waarvan de positieve
wortel is $S = [-\kappa + \sqrt{\kappa^2 + 4E^2x^2/S_c^2}\,]/(2Ex/S_c^2)$. Bij
20, 50, 100 en \qty{200}{m}: 0.36, 0.68, 0.89 en 1.03, tegenover 0.4, 1.0,
2.0 en 4.0 voor lineaire diffusie ($S = Ex/\kappa$). De heuvel heeft een
afgeronde top, daarna zijden waarvan de helling naar $S_c$ streeft: bijna
vlakke hellingen, met veel minder reliëf dan de parabool.
\end{solution}

\begin{solution}{exo:b3:quantitative-geomorphology:11}
Volgens de kinematische-golfstelling bereikt de verandering een punt met
coördinaat $\chi$ na $\chi/(KA_0^m)$; het netwerk heeft volledig gereageerd wanneer het
de laatste bronlocatie heeft bereikt, met de grootste $\chi$. $K A_0^m = 10^{-6}\times
1000 = \qty{e-3}{m/yr}$: de hoofdstroom reageert in \qty{5.46}{Myr}, het netwerk
in \qty{7.24}{Myr}, de orde van een miljoen jaar zoals verwacht.
\end{solution}

\begin{solution}{exo:b3:quantitative-geomorphology:12}
Boven de samenvloeiing is de stationaire helling van de zijrivier $S = (U/2K)A^{-\theta}$,
dus $\mathrm{d}z/\mathrm{d}\chi = (U/2K)A_0^{-\theta}$, de helft van die van de hoofdstroom.
Het \omterm{def:b3:quantitative-geomorphology:knickpoint}{knikpunt} bereikt de samenvloeiing wanneer dat van de hoofdstroom bij $\chi_j$ ligt; na een
verdere tijd $\Delta t$ is dat van de hoofdstroom verschoven naar $\chi_k$ met $\chi_k -
\chi_j = K A_0^m\,\Delta t$, terwijl dat van de zijrivier, met celeriteit $2KA^m$,
is naar $\chi$ verschoven met $\chi - \chi_j = 2KA_0^m\,\Delta t = 2(\chi_k -
\chi_j)$.
\end{solution}

\begin{solution}{pb:b3:quantitative-geomorphology:1}
\textbf{1.} $4.5\times12000^{1.7} = \qty{3.87e7}{m^2}$, \qty{38.7}{km^2};
$4.5\times300^{1.7} = \qty{7.3e4}{m^2}$, \qty{0.073}{km^2}.

\textbf{2.} $\theta \approx -\log(0.007855/0.06641)/\log(37.08/0.388) =
0.928/1.980 = 0.47$.

\textbf{3.} Ja: met een spreiding van ongeveer \qty{10}{\percent} op twaalf punten,
$0.466$ verschilt niet significant van 0.5.

\textbf{4.} $0.06641\times\sqrt{388\,000} = \qty{41.4}{m}$;
$0.007855\times\sqrt{3.708\times10^7} = \qty{47.8}{m}$. Het geometrisch gemiddelde,
\qty{45.1}{m}, schat $k_s$ voor $\theta = 0.5$.

\textbf{5.} Een helling is het verschil van twee hoogten, elk onzeker tot op
decimeters tot meters, gedeeld door een korte afstand; op zachte trajecten is de
fout een groot deel van het verschil. $\chi$ integreert in plaats van
te differentiëren.

\textbf{6.} $p = 1 - 1.7\times0.5 = 0.15$; $\chi = (1000/\sqrt{4.5})\,
(12000^{0.15} - 9000^{0.15})/0.15 = \qty{543.1}{m}$.

\textbf{7.} $(193.4 - 170)/543.1 = 0.0431$; $k_{sn} = 0.0431\times A_0^{0.5}
= \qty{43.1}{m}$.

\textbf{8.} De chi-schatting rust op hoogten over \qty{3}{km} rivierlengte,
niet op twaalf ruisachtige hellinggradiënten.

\textbf{9.} $K = U/k_{sn} = 4.3\times10^{-5}/43.1 = \qty{1.00e-6}{yr^{-1}}$.

\textbf{10.} $10^{-6}\times\sqrt{3.87\times10^7}\times0.00691 =
\qty{4.30e-5}{m/yr}$, \qty{43}{m/Myr}: de uitlaat houdt gelijke tred met de
opheffing.

\textbf{11.} Bij de uitlaat $10^{-6}\times6222 = \qty{6.2}{m/kyr}$; bij de
waterval, $d = \qty{7.40}{km}$, $A = \qty{1.70e7}{m^2}$ en
\qty{4.1}{m/kyr}.

\textbf{12.} $170 + 43\times(4.5)^{-0.5}(12000^{0.15} - 300^{0.15})/0.15 =
\qty{405}{m}$: een reliëf van \qty{235}{m}.

\textbf{13.} Voor $n = 1$ is de vergelijking een kinematische golf: punten aan beide
zijden van de trap stijgen met $U$ per tijdseenheid, zodat de trap zijn hoogte behoudt,
de \qty{25}{m} van de val.

\textbf{14.} $\chi = (1000/\sqrt{4.5})(12000^{0.15} - 7400^{0.15})/0.15 =
\qty{899.4}{m}$.

\textbf{15.} $t = \chi/(KA_0^{0.5}) = 899.4/10^{-3} = \qty{899}{kyr}$.

\textbf{16.} Bij $\chi = 899.4 + 300 = \qty{1199}{m}$:
$d^{0.15} = 12000^{0.15} - 1199\times0.15\times\sqrt{4.5}/1000$, $d =
\qty{6.25}{km}$, \qty{5.75}{km} stroomopwaarts van de uitlaat.

\textbf{17.} Bij de samenvloeiing ($d = \qty{9.5}{km}$) $\chi = \qty{442.8}{m}$;
de waterval ligt waar de $\chi$ van de zijrivier \qty{899.4}{m} bereikt:
$s^{0.15} = 5000^{0.15} - (899.4 - 442.8)\times0.15\times\sqrt{1.5}/1000$,
$s = \qty{4.27}{km}$, \qty{0.73}{km} boven de samenvloeiing.

\textbf{18.} Hun samenvloeiingen liggen bij $\chi = 1954$ en \qty{999}{m}, voorbij
\qty{899}{m}: het \omterm{def:b3:quantitative-geomorphology:knickpoint}{knikpunt} is er nog niet langs getrokken, dus de hoofdrivier is daar nog
op zijn oude niveau.

\textbf{19.} $E = 5\times10^{-3}\times0.0086 = \qty{4.3e-5}{m/yr}$,
\qty{43}{m/Myr}: de heuvels boven de waterval eroderen met de opheffingssnelheid, in
\omterm{def:b3:quantitative-geomorphology:steady-profile}{stationaire toestand}.

\textbf{20.} $150^2/(5\times10^{-3}) = \qty{4.5}{Myr}$, vijf keer de ouderdom van
de waterval.

\textbf{21.} Hun voet is met de geul \qty{25}{m} verlaagd: hun
onderste delen zijn te steil geworden en glijden misschien, terwijl hun kammen nog
met de oude snelheid eroderen. Een golf van steiler worden klimt elke helling op; na enkele
miljoen jaar keren ze terug naar dezelfde parabool, \qty{25}{m} lager.

\textbf{22.} $t = \chi k_{sn}/(U A_0^{0.5})$ varieert als $1/U$: van
$899/1.2 = 749$ tot $899/0.8 = \qty{1124}{kyr}$.

\textbf{23.} Bij $n = 2$ groeit de celeriteit met de helling: het steile
vlak zou sneller bewegen dan de trajecten eromheen en scherp blijven, maar
zijn snelheid en de daaruit af te lezen ouderdom zouden verschillen.

\textbf{24.} Het \omterm{def:b3:quantitative-geomorphology:knickpoint}{basisniveau} van de Varden-rivier daalde
\textbf{\qty[text-series-to-math=true]{900}{kyr}} geleden (\qty{899}{kyr} volgens de
berekening).
\end{solution}
